\section{Related work}
\label{sec:related}

\paragraph{RL for biological sequence design.} The field's 2025 main-track methods include $\delta$-Conservative Search \citep{kim2025deltacs}, which adapts the exploration radius using proxy uncertainty as a hedge against out-of-distribution reward inflation, $\mu$Protein and $\mu$Search \citep{sun2025muprotein}, which fit an explicit epistasis-aware predictor and run multi-step RL on top, and the Caltech / Yue-lab steering work of \citet{yang2025steering}, which conditions a protein language model on a desired-fitness trajectory rather than a scalar target. Bayesian active learning is represented by ALDE \citep{yang2025alde}; gradient-free RL by GLID$^2$E \citep{cao2025glid2e}; DNA-LLM RLVR by BioReason \citep{bioreason2025}. Earlier diversity-aware approaches include GFlowNets \citep{jain2022biological}, AdaLead \citep{sinai2020adalead}, PEX \citep{ren2022proximal}, and DyNA-PPO \citep{angermueller2020dyna}; latent-space RL for rugged proteins is treated by LatProtRL \citep{lee2024latprotrl}, and smoothed-landscape RL by \citet{kirjner2024smoothed}. \emph{None} of these methods decomposes the policy gradient into selection and transmission components. All of them inherit the undifferentiated-gradient pathology described in §\ref{sec:introduction}: the policy gradient pools the two effects, so the algorithm cannot tell, per round, whether it is exploiting a known peak or sampling an unexplored one. Our contribution is orthogonal to these methods rather than competitive with them: PRICE-RL's decomposition could in principle wrap any of them.

\paragraph{The Price equation in evolutionary biology and learning theory.} \citet{price1970selection} gave the exact decomposition $\Delta\bar f = \mathrm{Cov}(w,f)/\bar w + \mathbb{E}[w \cdot \Delta f]/\bar w$, relating the change in mean trait value across an evolving population to a selection covariance and a transmission expectation, and the decomposition has been used as an analytical tool in evolutionary biology for fifty years (see \citealp{frank2012natural} for a review). Its connection to learning algorithms remained essentially metaphorical until \citet{frank2025price} proved that the algebraic identity $\Delta\theta = M f + b + \xi$ specialises to natural selection, Bayesian updating, Newton's method, SGD, Adam, and the policy gradient as instances of the same force--metric--bias law. That paper supplied the unifying algebra; it did not propose any algorithm, controller, or diagnostic that exploited the decomposition as an active component of an update rule. PRICE-RL is, to our knowledge, the first algorithmic instantiation, and it is the autocorrelation-parameterised controller target plus the per-round Price ratio that turn the algebra into a piece of working machinery rather than a re-description of one.

\paragraph{Reward hacking detection.} \citet{skalse2022defining} formalised reward hacking as the gap between proxy and true reward under policy improvement, and \citet{laidlaw2025correlated} introduced occupancy-measure mitigations that bound how far a policy may drift from a reference distribution. Neither of these connects to a per-update diagnostic; both are post-hoc tests on aggregate behaviour. The proxy-uncertainty signal used by $\delta$-CS \citep{kim2025deltacs} is the closest comparator and operates round-by-round, but it lives at the predictor side rather than the gradient side: it asks \emph{is the surrogate's prediction reliable on this batch?} The Price ratio asks a different and earlier question: \emph{is the gradient consuming new sequence space or refining what the policy already supports?} Reward hacking is structurally a refinement-on-artefact event, and so $\rho_t$ flags it without waiting for variance to accumulate. We find empirically that the lead is $9$--$18$ rounds on $100\%$ of seeds, and the diagnostic is robust to surrogate corruption up to $\sigma_n = 0.40$.

\paragraph{Two-timescale and structured policy gradients.} Methods such as actor--critic \citep{konda2000actor} use different rates for different function approximators (policy versus value); two-timescale variance-reduction analyses \citep{karimi2022two} use them for distinct stochastic-approximation processes. PRICE-RL uses different rates for \emph{decomposed components of the same policy gradient}, which is a different operation: the two timescales live inside one estimator, partitioned by support membership rather than by network identity. This is why the decomposition's exactness theorem (Theorem~\ref{thm:exact}) is sample-wise rather than asymptotic, and why a controller targeting a single scalar $\rho_t$ is sufficient to reproduce the regret-optimal mixing in our experiments. The structured policy-gradient literature (e.g., advantage decompositions, separated value heads) addresses variance; the Price decomposition addresses interpretability and steering of the gradient itself.
